\section{Light --- Lib --- Count Sort}
%%% generated by the blueprint pipeline from TCSlib.Programs.Light.Lib.CountSort
\subsection*{Light --- Lib --- Count Sort: Introduction}

Declarations of the Lean module \texttt{TCSlib.Programs.Light.Lib.CountSort}.

\subsection*{Light --- Lib --- Count Sort: Results}

\begin{definition}[Length index]
\lean{Light.CountSort.Len}
\label{Light.CountSort.Len}
\leanok
The length index is the natural number $0$.
\end{definition}

\begin{definition}[Source index]
\lean{Light.CountSort.Src}
\label{Light.CountSort.Src}
\leanok
\uses{Light.CountSort.Len}
The source index is the natural number $1$.
\end{definition}

\begin{definition}[Destination index]
\lean{Light.CountSort.Dst}
\label{Light.CountSort.Dst}
\leanok
\uses{Light.CountSort.Len}
The destination index is the natural number $2$.
\end{definition}

\begin{definition}[Keys index]
\lean{Light.CountSort.Keys}
\label{Light.CountSort.Keys}
\leanok
\uses{Light.CountSort.Len}
The keys index is the natural number $3$.
\end{definition}

\begin{definition}[Stride index]
\lean{Light.CountSort.Stride}
\label{Light.CountSort.Stride}
\leanok
\uses{Light.CountSort.Len}
The stride index is the natural number $4$.
\end{definition}

\begin{definition}[Offset constant]
\lean{Light.CountSort.Offset}
\label{Light.CountSort.Offset}
\leanok
\uses{Light.CountSort.Len}
The offset is the natural number $5$.
\end{definition}

\begin{definition}[Bucket count]
\lean{Light.CountSort.Buckets}
\label{Light.CountSort.Buckets}
\leanok
\uses{Light.CountSort.Len}
The number of buckets is the natural number $6$.
\end{definition}

\begin{definition}[Count index]
\lean{Light.CountSort.Cnt}
\label{Light.CountSort.Cnt}
\leanok
\uses{Light.CountSort.Len}
The count index is the natural number $7$.
\end{definition}

\begin{definition}[Index constant]
\lean{Light.CountSort.Idx}
\label{Light.CountSort.Idx}
\leanok
\uses{Light.CountSort.Len}
The index is the natural number $8$.
\end{definition}

\begin{definition}[Cell constant]
\lean{Light.CountSort.Cell}
\label{Light.CountSort.Cell}
\leanok
\uses{Light.CountSort.Len}
The cell is the natural number $9$.
\end{definition}

\begin{definition}[Place constant]
\lean{Light.CountSort.Place}
\label{Light.CountSort.Place}
\leanok
\uses{Light.CountSort.Len}
The place is the natural number $10$.
\end{definition}

\begin{definition}[Counting-sort key expression]
\lean{Light.csKey}
\label{Light.csKey}
\leanok
\uses{Light.CountSort.Idx, Light.CountSort.Keys, Light.CountSort.Offset, Light.CountSort.Src, Light.CountSort.Stride, Light.Expr, Light.M, Light.v}
The key expression is the memory access at the address obtained by adding the offset to
the product of the stride and the memory access at the sum of the source and index
variables, then adding the keys variable.
\end{definition}

\begin{theorem}[Cost of the key expression]
\lean{Light.csKey_cost}
\label{Light.csKey_cost}
\leanok
\uses{Light.Expr.cost, Light.csKey}
\difficulty{1}
The expression used to load an item's key in counting sort has cost $11$.
\end{theorem}

\begin{definition}[Zeroing the bucket counts]
\lean{Light.csZero}
\label{Light.csZero}
\leanok
\uses{Light.CountSort.Buckets, Light.CountSort.Cnt, Light.CountSort.Idx, Light.Stmt, Light.Stmt.for, Light.k, Light.v}
The command iterates over the indices from $0$ through the bucket count and stores $0$
at each address equal to the count variable plus the current index.
\end{definition}

\begin{definition}[Counting Sort Frequency Loop]
\lean{Light.csCount}
\label{Light.csCount}
\leanok
\uses{Light.CountSort.Cell, Light.CountSort.Cnt, Light.CountSort.Idx, Light.CountSort.Len, Light.M, Light.Stmt, Light.Stmt.for, Light.csKey, Light.k, Light.v}
The frequency-counting step of counting sort is the loop over all $n$ input items that
computes the item's key-cell address as the count-array base plus the key plus one, then
increments the value stored in that cell by one.
\end{definition}

\begin{definition}[Counting Sort Prefix-Sum Loop]
\lean{Light.csPrefix}
\label{Light.csPrefix}
\leanok
\uses{Light.CountSort.Buckets, Light.CountSort.Cnt, Light.CountSort.Idx, Light.M, Light.Stmt, Light.Stmt.for, Light.k, Light.v}
The prefix-sum step of counting sort is the loop over all bucket indices that replaces
the count at each index plus one with the sum of the counts at that index and the next
index.
\end{definition}

\begin{definition}[Counting Sort Placement Loop]
\lean{Light.csPlace}
\label{Light.csPlace}
\leanok
\uses{Light.CountSort.Cell, Light.CountSort.Cnt, Light.CountSort.Dst, Light.CountSort.Idx, Light.CountSort.Len, Light.CountSort.Place, Light.CountSort.Src, Light.M, Light.Stmt, Light.Stmt.for, Light.csKey, Light.k, Light.v}
The placement step of counting sort loops over all $n$ input items. For each item, it
reads the current placement value from the count cell indexed by the item's key, writes
the item from the source array to the destination array at that position, and increments
the placement value by one.
\end{definition}

\begin{definition}[Counting sort count-array shift]
\lean{Light.csShift}
\label{Light.csShift}
\leanok
\uses{Light.CountSort.Buckets, Light.CountSort.Cnt, Light.CountSort.Idx, Light.M, Light.Stmt, Light.k, Light.v}
Let $	ext{Idx}$ and $	ext{Cnt}$ be program variables, let $	ext{Buckets}$ denote the
bucket count, and let $M(a)$ denote the value read from memory address $a$. The program
sets $	ext{Idx}$ to $	ext{Buckets}$; while $	ext{Idx}>0$, it copies the value at address
$	ext{Cnt}+	ext{Idx}-1$ to address $	ext{Cnt}+	ext{Idx}$ and decrements $	ext{Idx}$;
then it writes $0$ to address $	ext{Cnt}$.
\end{definition}

\begin{definition}[Counting Sort Body]
\lean{Light.countSortBody}
\label{Light.countSortBody}
\leanok
\uses{Light.Stmt, Light.csCount, Light.csPlace, Light.csPrefix, Light.csShift, Light.csZero}
The body of counting sort executes, in order, the zero-initialization step, the
frequency-counting step, the prefix-sum step, the placement step, and the count-array
shift step.
\end{definition}

\begin{definition}[Counting Sort Arguments]
\lean{Light.CountSort.Args}
\label{Light.CountSort.Args}
\leanok
The counting-sort arguments consist of a number $n$ of items; source and destination
locations; a key-field base, stride, and offset specifying where each item's key is
stored; a bound $nb$ such that all keys are less than $nb$; and a base location $cnt$
for $nb+1$ scratch cells.
\end{definition}

\begin{definition}[Key Cell Address]
\lean{Light.CountSort.Args.keyCell}
\label{Light.CountSort.Args.keyCell}
\leanok
\uses{Light.CountSort.Args}
For counting-sort arguments with key-field base $key$, stride $stride$, and offset
$off$, the address of the key cell for item $x$ is $key+x\cdot stride+off$.
\end{definition}

\begin{definition}[Counting Sort State Construction]
\lean{Light.CountSort.Args.state}
\label{Light.CountSort.Args.state}
\leanok
\uses{Light.CountSort.Args, Light.State, Light.frame}
For counting-sort arguments $A$, integers $i,c,p$, and a memory function $\mu:\mathbb{N}\to\mathbb{Z}$,
$A.\mathrm{state}(i,c,p,\mu)$ is the light state whose local frame contains, in order,
$A.n,A.src,A.dst,A.key,A.stride,A.off,A.nb,A.cnt,i,c,p$, and whose memory is $\mu$.
\end{definition}

\begin{definition}[Agreement Outside Output and Scratch Regions]
\lean{Light.CountSort.Args.Same}
\label{Light.CountSort.Args.Same}
\leanok
\uses{Light.CountSort.Args, Light.SameOutside2}
For counting-sort arguments $A$, the relation $A.\mathrm{Same}(\mu,\mu\prime)$ means that
the two memories agree at every address outside both the output interval of length $A.n$
starting at $A.dst$ and the scratch interval of length $A.nb+1$ starting at $A.cnt$.
\end{definition}

\begin{definition}[Counting Sort Preconditions]
\lean{Light.CountSort.Pre}
\label{Light.CountSort.Pre}
\leanok
\uses{Light.Apart, Light.CountSort.Args, Light.CountSort.Args.keyCell, Light.Limits, Light.Outside}
For limits $lim$, memory $\mu:\mathbb{N}\to\mathbb{Z}$, counting-sort arguments $A$, and
functions $it,kf:\mathbb{N}\to\mathbb{N}$, the precondition requires
$lim.space\le lim.word$, and that the source interval $[A.src,A.src+A.n)$, destination
interval $[A.dst,A.dst+A.n)$, and counter interval $[A.cnt,A.cnt+A.nb+1)$ all fit within
the space bound and are pairwise disjoint. For every $i<A.n$, the source cell contains
$it(i)$; the key address $A.keyCell(it(i))$ is below $lim.space$, lies outside both the
destination and counter intervals, and contains $kf(i)$; and $kf(i)<A.nb$.
\end{definition}

\begin{definition}[Counting sort postconditions]
\lean{Light.CountSort.Post}
\label{Light.CountSort.Post}
\leanok
\uses{Light.CountSort.Args, Light.CountSort.Args.Same, Light.cntLt, Light.sortPos}
For counting-sort arguments $A$, initial memory $\mu$, item and key functions $it,kf$, and
final memory $\mu\prime$, the postcondition says that every item $j<A.n$ is stored at
destination address $A.dst+sortPos(kf,A.n,j)$ with value $it(j)$; for each $t\le A.nb$,
the counter cell $A.cnt+t$ contains the number of items among the first $A.n$ whose key
is less than $t$; and $\mu$ and $\mu\prime$ agree outside the destination and counter
intervals specified by $A.\mathrm{Same}$.
\end{definition}

\begin{theorem}[Preservation of source items]
\lean{Light.CountSort.Pre.item}
\label{Light.CountSort.Pre.item}
\leanok
\uses{Light.CountSort.Args.Same, Light.CountSort.Pre}
\difficulty{3}
Let $\mu,\mu\prime:\mathbb{N}\to\mathbb{Z}$ be memories, $A$ counting-sort arguments, and
$it,kf:\mathbb{N}\to\mathbb{N}$ item and key functions. If $\mu$ satisfies the counting-sort
precondition, $\mu\prime$ agrees with $\mu$ outside the destination and counter regions,
and $i<A.n$, then $\mu\prime(A.src+i)=it(i)$.
\end{theorem}

\begin{theorem}[Key expression evaluates safely]
\lean{Light.csKey_gives}
\label{Light.csKey_gives}
\leanok
\uses{Light.CountSort.Args.Same, Light.CountSort.Args.keyCell, Light.CountSort.Args.state, Light.CountSort.Pre, Light.CountSort.Pre.item, Light.Expr.Gives, Light.Limits.Addr, Light.csKey}
\difficulty{5}
Let $lim$ be resource limits, $\mu,\mu\prime$ memories, $A$ counting-sort arguments, and
$it,kf:\mathbb{N}\to\mathbb{N}$ item and key functions. Suppose the precondition holds
for $lim,\mu,A,it,kf$, the two memories agree outside the destination and counter
regions, and $i<A.n$. For any integers $c,p$, the key-reading expression is safe in the
state $A.\mathrm{state}(i,c,p,\mu\prime)$ under $lim$ and evaluates to $kf(i)$.
\end{theorem}

\begin{definition}[Counting Sort State Predicate]
\lean{Light.CsState}
\label{Light.CsState}
\leanok
\uses{Light.CountSort.Args, Light.CountSort.Args.Same, Light.CountSort.Args.state, Light.State}
Let $\mu$ be an initial memory, $A$ the counting sort parameters, $i,c,p\in\mathbb{Z}$,
$R$ a predicate on memories, and $\sigma$ a state. The state satisfies the counting sort
state predicate if there are integers $c,p$ and a memory map $\mu'$ such that $\sigma$
is the state associated with $A,i,c,p,\mu'$, the memories $\mu$ and $\mu'$ agree outside
the destination and counter regions, and $R(\mu')$ holds.
\end{definition}

\begin{theorem}[Implication for count-sort states]
\lean{Light.CsState.imp}
\label{Light.CsState.imp}
\leanok
\uses{Light.CsState}
\difficulty{4}
Let $\mu$ be an initial memory, $A$ counting-sort arguments, $e\in\mathbb{Z}$, $R,R'$
predicates on integer-valued memories, and $\sigma$ a state. If $\sigma$ can be
represented using entry value $e$ and a memory $\mu'$ that agrees with $\mu$ outside the
destination and scratch regions and satisfies $R$, and every memory satisfying $R$ also
satisfies $R'$, then $\sigma$ has such a representation with a memory satisfying $R'$.
\end{theorem}

\begin{theorem}[Counting-sort loop with an invariant]
\lean{Light.Ends.csFor}
\label{Light.Ends.csFor}
\leanok
\uses{Light.CountSort.Args.Same, Light.CountSort.Args.state, Light.CountSort.Idx, Light.CsState, Light.Ends, Light.Ends.forShape, Light.Ends.mono, Light.Expr, Light.Expr.Gives, Light.Expr.cost, Light.Stmt, Light.Stmt.blockCost, Light.Stmt.for}
\difficulty{5}
Let $L$ be resource limits, $P$ a program, and $d$ an execution depth. Let $A$ be
counting-sort arguments, let $\mu:\mathbb{N}\to\mathbb{Z}$ be the original memory, let
$\sigma$ be an initial state, let $e\in\mathbb{Z}$ be an entry value, let $h$ be an
upper-bound expression, let $s$ be a loop body, let $T,m\in\mathbb{N}$, and let $R$ be a
predicate on indices and memories. Suppose $\sigma$ satisfies the counting-sort state
condition for $A$ at entry value $e$, with invariant $R(0)$. For every $j<m$, all
integer counters $c,p$, and every memory $\mu'$ that agrees with $\mu$ outside both the
destination interval and the scratch interval and satisfies $R(j)$, suppose the body
executes from the state encoded by $A,j,c,p,\mu'$ within the body's block cost and ends
in a state satisfying the counting-sort state condition at entry value $j$ with
invariant $R(j+1)$. Suppose also that, for every $j\in\mathbb{N}$, all integers $c,p$,
and every memory $\mu'$, the expression $h$ is safe in the state encoded by
$A,j,c,p,\mu'$ and has value $m$ there; that $m\leq L.word$; and that
$m(\operatorname{cost}(h)+\operatorname{blockCost}(s)+7)+\operatorname{cost}(h)+5\leq
T$. Then the for-loop with counter variable $8$, upper bound $h$, and body $s$ executes
from $\sigma$ within cost $T$ and ends in a state satisfying the counting-sort state
condition for $A$ at entry value $m$ with invariant $R(m)$.
\end{theorem}

\begin{definition}[Zeroed Counter Prefix]
\lean{Light.CsZeroed}
\label{Light.CsZeroed}
\leanok
\uses{Light.CountSort.Args}
Let $A$ be a counting-sort argument, let $j\in\mathbb{N}$, and let $\mu'$ be an
integer-valued memory map. The first $j$ counter cells, starting at the counter region's
initial address, contain zero: for every $t<j$, the memory value at address $A.cnt+t$ is
zero.
\end{definition}

\begin{theorem}[Extension of a zeroed counter prefix]
\lean{Light.CsZeroed.succ}
\label{Light.CsZeroed.succ}
\leanok
\uses{Light.CsZeroed}
\difficulty{4}
Let $A$ be counting-sort arguments, $j\in\mathbb{N}$, and $\mu'$ an integer-valued
memory map. If every counter at offset $t<j$ contains zero, then changing the counter at
offset $j$ to zero makes every counter at offset $t<j+1$ contain zero.
\end{theorem}

\begin{theorem}[Counting-sort counter initialization]
\lean{Light.csZero_ends}
\label{Light.csZero_ends}
\leanok
\uses{Light.CountSort.Pre, Light.CsState, Light.CsState.imp, Light.CsZeroed, Light.CsZeroed.succ, Light.Ends, Light.Ends.csFor, Light.SameOutside2.update, Light.csZero}
\difficulty{4}
Let a counting-sort input have $n$ items and $N$ buckets, with the word and address
bounds and required disjointness conditions for the source, destination, and counter
ranges and key cells. From any state with arbitrary loop index, other loop variables,
and memory agreeing with the initial memory outside the destination and counter ranges,
the counter-initialization program has an execution of cost at most $15N+23$ ending with
loop index $N+1$, with counters $0$ through $N$ all equal to zero and memory still
agreeing with its initial value outside the destination and counter ranges.
\end{theorem}

\begin{definition}[Counted key frequencies]
\lean{Light.CsCounted}
\label{Light.CsCounted}
\leanok
\uses{Light.CountSort.Args, Light.cntEq}
For arguments $A$, a key function $k:\mathbb{N}\to\mathbb{N}$, a count $i\in\mathbb{N}$, and
memory $\mu\prime:\mathbb{N}\to\mathbb{Z}$, the counted-key condition says that
$\mu\prime(A.cnt)=0$ and, for every key $t<A.nb$, the cell $A.cnt+(t+1)$ contains
the number of indices $x<i$ for which $k(x)=t$.
\end{definition}

\begin{theorem}[Initial counter state]
\lean{Light.CsCounted.start}
\label{Light.CsCounted.start}
\leanok
\uses{Light.CsCounted, Light.CsZeroed}
\difficulty{3}
Let $A$ be counting-sort arguments, $kf:\mathbb{N}\to\mathbb{N}$ a key function, and
$\mu\prime:\mathbb{N}\to\mathbb{Z}$ a memory. If the first $A.nb+1$ counter cells
are zero, then the initial counter state records that no items have been counted: the
cell $A.cnt$ is zero and, for each key $t<A.nb$, the cell $A.cnt+(t+1)$ contains the
number of indices $i<0$ with $kf(i)=t$ (which is zero).
\end{theorem}

\begin{theorem}[Successor update of counted frequencies]
\lean{Light.CsCounted.succ}
\label{Light.CsCounted.succ}
\leanok
\uses{Light.CsCounted, Light.cntEq, Light.cntEq_succ}
\difficulty{4}
Let $A$ be counting-sort arguments, let $kf\colon\mathbb{N}\to\mathbb{N}$ assign a key
to each index, and let $i\in\mathbb{N}$. Suppose an integer-valued memory map $\mu'$ has
value $0$ at address $A.cnt$, and for every key $t<A.nb$ has at address $A.cnt+(t+1)$
the number of indices $j<i$ with $kf(j)=t$. Change the value at address
$A.cnt+(kf(i)+1)$ to one more than the number of indices $j<i$ with $kf(j)=kf(i)$,
leaving every other address unchanged. The resulting memory map has value $0$ at
$A.cnt$, and at each address $A.cnt+(t+1)$ for $t<A.nb$ it stores the number of indices
$j<i+1$ with $kf(j)=t$.
\end{theorem}

\begin{theorem}[Counting-sort counting-loop execution]
\lean{Light.csCount_ends}
\label{Light.csCount_ends}
\leanok
\uses{Light.CountSort.Pre, Light.CsCounted, Light.CsCounted.start, Light.CsCounted.succ, Light.CsState, Light.CsState.imp, Light.CsZeroed, Light.Ends, Light.Ends.csFor, Light.SameOutside2.update, Light.cntEq, Light.csCount, Light.csKey_gives}
\difficulty{5}
Let $P$ be a program, let $d\in\mathbb{N}$, let $\mathrm{lim}$ specify word, space, and
depth limits, and take counting-sort parameters $A$, initial memory $\mu$, item and key
functions $\mathrm{it},kf:\mathbb{N}\to\mathbb{N}$, and an initial state $\sigma$.
Suppose the counting-sort preconditions hold. Suppose further that for some integer
$\mathrm{entry}$, $\sigma$ is the state formed from the standard counting-sort argument
frame with index $\mathrm{entry}$, arbitrary integer counter and place variables, and a
memory $\mu'$ that agrees with $\mu$ outside the destination and counter intervals;
moreover, every counter cell $\mathrm{cnt}+t$ for $t<\mathrm{nb}+1$ contains zero. Then
the counting loop has an execution from $\sigma$ of cost at most $30n+6$ whose final
state has index $n$, still agrees with $\mu$ outside those intervals, has counter cell
$\mathrm{cnt}$ equal to zero, and has cell $\mathrm{cnt}+(t+1)$ equal to the number of
indices $i<n$ with $kf(i)=t$ for every $t<\mathrm{nb}$.
\end{theorem}

\begin{definition}[Summed key counts]
\lean{Light.CsSummed}
\label{Light.CsSummed}
\leanok
\uses{Light.CountSort.Args, Light.cntEq, Light.cntLt}
Let $A$ specify the number of items and the counter region, let $kf$ assign a key to
each item index, let $r$ be a natural number, and let $\mu'$ be a memory map from
natural-number addresses to integers. For every natural number $t\le r$, the cell at
address $A.cnt+t$ contains the number of indices $i<A.n$ whose key is less than $t$. For
every natural number $t<A.nb$ with $r\le t$, the cell at address $A.cnt+(t+1)$ contains
the number of indices $i<A.n$ whose key equals $t$.
\end{definition}

\begin{theorem}[Initial summed-counter invariant]
\lean{Light.CsSummed.start}
\label{Light.CsSummed.start}
\leanok
\uses{Light.CsCounted, Light.CsSummed}
\difficulty{4}
Let $A$ be counting-sort arguments, let $kf:\mathbb{N}\to\mathbb{N}$ assign keys to item
indices, and let $\mu'$ be an integer-valued memory map. Suppose the cell at address
$A.cnt$ contains $0$, and for every key $t<A.nb$, the cell at address $A.cnt+(t+1)$
contains the number of indices $i<A.n$ with $kf(i)=t$. Then for every $t\le 0$, the cell
at address $A.cnt+t$ contains the number of indices $i<A.n$ with $kf(i)<t$, and for
every $t<A.nb$, the cell at address $A.cnt+(t+1)$ contains the number of indices $i<A.n$
with $kf(i)=t$.
\end{theorem}

\begin{theorem}[Successor summed-counter invariant]
\lean{Light.CsSummed.succ}
\label{Light.CsSummed.succ}
\leanok
\uses{Light.CsSummed, Light.cntLt}
\difficulty{4}
Let $A$ be counting-sort arguments, $kf$ a natural-number key function, $r$ a natural
number, and $\mu\prime$ a memory satisfying the stage-$r$ counter conditions: for each
$t\le r$, cell $A.cnt+t$ contains the number of items among the first $A.n$ with key
less than $t$; and for each $r\le t<A.nb$, cell $A.cnt+(t+1)$ contains the number of
items with key exactly $t$. Updating cell $A.cnt+(r+1)$ to the number of items with key
less than $r+1$ gives the stage-$(r+1)$ counter conditions.
\end{theorem}

\begin{theorem}[Counting-sort prefix phase]
\lean{Light.csPrefix_ends}
\label{Light.csPrefix_ends}
\leanok
\uses{Light.CountSort.Pre, Light.CsCounted, Light.CsState, Light.CsState.imp, Light.CsSummed, Light.CsSummed.start, Light.CsSummed.succ, Light.Ends, Light.Ends.csFor, Light.SameOutside2.update, Light.cntLt, Light.cntLt_le, Light.cntLt_succ_left, Light.csPrefix}
\difficulty{5}
Let a counting-sort input have $n$ items with keys in $\{0,\ldots,N-1\}$, and assume the
address bounds, key-cell bounds, and required disjointness conditions for the source,
destination, and counter ranges. If the counters initially contain the frequencies of
keys among the first $n$ items (with counter zero equal to $0$), then the prefix-sum
program has an execution of cost at most $25N+6$ that ends with loop index $N$, counters
containing the cumulative frequencies through each key, and memory agreeing with its
initial value outside the destination and counter ranges.
\end{theorem}

\begin{definition}[Placement in Counting Sort]
\lean{Light.CsPlaced}
\label{Light.CsPlaced}
\leanok
\uses{Light.CountSort.Args, Light.cntEq, Light.cntLt, Light.sortPos}
Let $n$, $d_0$, $b$, and $c$ be natural numbers, where $n$ is the number of items, $d_0$
is the base address of the destination region, $b$ is the key bound, and $c$ is the base
address of the counter region. Let $k(j)$ be the key of item $j$, let $x(j)$ specify the
item at index $j$, and let $m$ be a memory map. For every $t<b$, the counter cell at
address $c+t$ contains the number of indices $j<n$ with $k(j)<t$, plus the number of
indices $j<i$ with $k(j)=t$. For every $j<i$, the destination cell at address
$d_0+	ext{sortPos}(k,n,j)$ contains $x(j)$, where $	ext{sortPos}(k,n,j)$ denotes the
sorted position used for item $j$ among the first $n$ items.
\end{definition}

\begin{theorem}[Initial Counting-Sort Placement State]
\lean{Light.CsPlaced.start}
\label{Light.CsPlaced.start}
\leanok
\uses{Light.CsPlaced, Light.CsSummed}
\difficulty{3}
Let $A$ be counting-sort arguments, let $kf$ assign a key to each item index, and let
$\mu'$ be an integer-valued memory map. Suppose that for every $t\le A.nb$, the cell at
address $A.cnt+t$ contains the number of indices $i<A.n$ whose key is less than $t$, and
that for every $t<A.nb$ with $A.nb\le t$, the cell at address $A.cnt+(t+1)$ contains the
number of indices $i<A.n$ whose key equals $t$. Then for every $t<A.nb$, the cell at
address $A.cnt+t$ contains the number of indices $i<A.n$ whose key is less than $t$, and
for every $j<0$, the destination cell at the stable sorting position of item $j$
contains that item.
\end{theorem}

\begin{theorem}[Successor placement invariant]
\lean{Light.CsPlaced.succ}
\label{Light.CsPlaced.succ}
\leanok
\uses{Light.CountSort.Pre, Light.CsPlaced, Light.cntEq_succ, Light.eq_of_sortPos_eq, Light.sortPos, Light.sortPos_lt}
\difficulty{5}
Let $n$, $src$, $dst$, $key$, $stride$, $off$, $nb$, and $cnt$ be natural numbers, and
let $it,kf:\mathbb{N}\to\mathbb{N}$ and $\mu,\mu':\mathbb{N}\to\mathbb{Z}$. Let the
resource limits be $word\in\mathbb{Z}$ and $space,depth\in\mathbb{N}$. Assume $space\le
word$; the source, destination, and counter intervals, of lengths $n,n,nb+1$
respectively, fit within $space$ and are pairwise apart; and for every $r<n$, the source
cell at $src+r$ contains $it(r)$, the key cell at $key+it(r)\cdot stride+off$ lies
within $space$ and outside both the destination and counter intervals, and contains
$kf(r)<nb$. Assume also that $i<n$ and that, for every $t<nb$, the cell at $cnt+t$ in
$\mu'$ contains the number of indices $r<n$ with $kf(r)<t$ plus the number of indices
$r<i$ with $kf(r)=t$, while for every $j<i$, the cell at $dst+p(j)$ contains $it(j)$,
where $p(j)$ is the number of indices $r<n$ with $kf(r)<kf(j)$ plus the number of
indices $r<j$ with $kf(r)=kf(j)$. Form a new memory map by first setting the cell at
$dst+p(i)$ to $it(i)$ and then setting the cell at $cnt+kf(i)$ to $p(i)+1$. In the
resulting map, for every $t<nb$ the cell at $cnt+t$ contains the number of indices $r<n$
with $kf(r)<t$ plus the number of indices $r<i+1$ with $kf(r)=t$, and for every $j<i+1$
the cell at $dst+p(j)$ contains $it(j)$.
\end{theorem}

\begin{theorem}[Counting-sort placement phase]
\lean{Light.csPlace_ends}
\label{Light.csPlace_ends}
\leanok
\uses{Light.CountSort.Pre, Light.CountSort.Pre.item, Light.CsPlaced, Light.CsPlaced.start, Light.CsPlaced.succ, Light.CsState, Light.CsState.imp, Light.CsSummed, Light.Ends, Light.Ends.csFor, Light.SameOutside2.update, Light.csKey_gives, Light.csPlace, Light.sortPos, Light.sortPos_lt}
\difficulty{5}
Let a counting-sort input have $n$ items $x_i$ with keys $k_i<N$, stored at source
address $s$, with a destination range and $N+1$ counter cells satisfying the required
bounds and disjointness conditions; assume the key cells are also within the address
bound and disjoint from both ranges. Suppose the counter memory has cumulative counts
$\#\{i<n:k_i<t\}$ at counter $t$ for $t\le N$, and bucket sizes $\#\{i<n:k_i=t\}$ at
counter $t+1$ for $t<N$. From any state with loop index $r$, arbitrary other loop
variables, and memory agreeing with the initial memory outside the destination and
counter ranges, the placement program has an execution of cost at most $38n+6$ ending
with loop index $n$, memory still agreeing outside those ranges, updated counters
encoding the next positions, and each processed item $j<n$ stored at its stable sorted
position.
\end{theorem}

\begin{definition}[Shifted counting-sort state]
\lean{Light.CsShifted}
\label{Light.CsShifted}
\leanok
\uses{Light.CountSort.Args, Light.cntLt, Light.sortPos}
For arguments $A$, item and key functions $it,kf:\mathbb{N}\to\mathbb{N}$, an index $s$, and
memory $\mu\prime$, the shifted-state condition requires that each item $j<A.n$ is at
destination address $A.dst+sortPos(kf,A.n,j)$ with value $it(j)$; for every $t\le A.nb$
with $s<t$, counter cell $A.cnt+t$ contains $cntLt(kf,t,A.n)$; and for every $t<s$,
counter cell $A.cnt+t$ contains $cntLt(kf,t+1,A.n)$.
\end{definition}

\begin{theorem}[Initial shifted-state invariant]
\lean{Light.CsShifted.start}
\label{Light.CsShifted.start}
\leanok
\uses{Light.CsPlaced, Light.CsShifted, Light.cntLt_succ_left}
\difficulty{4}
Let $A$ specify the number $n$ of items, the key bound $nb$, and the destination and
counter regions, and let $it$ and $kf$ give the item and key at each index. If every
item $j<n$ is stored at its stable sorting position—the number of indices $i<n$ with
$kf(i)<kf(j)$ plus the number of indices $i<j$ with $kf(i)=kf(j)$—then every item
remains at that position, and for every $t\le nb$ the counter cell at address $A.cnt+t$
contains the number of indices $i<n$ with $kf(i)<t$.
\end{theorem}

\begin{theorem}[Placement survives a counter update]
\lean{Light.CsShifted.placed_update}
\label{Light.CsShifted.placed_update}
\leanok
\uses{Light.CountSort.Pre, Light.CsShifted, Light.sortPos, Light.sortPos_lt}
\difficulty{3}
Let $A$ be counting-sort arguments, let $it,kf:\mathbb{N}\to\mathbb{N}$, and let $s$ be
a natural number. Suppose a memory map $\mu'$ places each item $j<A.n$ at address
$A.dst+p_j$, where $p_j$ is the number of indices $i<A.n$ with $kf(i)<kf(j)$ plus the
number of indices $i<j$ with $kf(i)=kf(j)$. For every $t\le A.nb$, every integer $x$,
and every $j<A.n$, changing the memory value at address $A.cnt+t$ to $x$ leaves the
value at address $A.dst+p_j$ equal to $it(j)$.
\end{theorem}

\begin{theorem}[Predecessor shifted counting-sort invariant]
\lean{Light.CsShifted.succ}
\label{Light.CsShifted.succ}
\leanok
\uses{Light.CountSort.Pre, Light.CsShifted, Light.CsShifted.placed_update, Light.cntLt}
\difficulty{4}
Let $A$ describe $n$ items, source and destination regions, key-cell addresses, a key
bound $N_b$, and a counter region; let $it,kf:\mathbb{N}\to\mathbb{N}$ and let
$\mu,\mu':\mathbb{N}\to\mathbb{Z}$ be memory maps. Assume the space bound is at most the
word bound, the source, destination, and counter regions fit in memory and are pairwise
disjoint, and for every $i<n$ the source contains $it(i)$, its key cell lies in memory
and outside the destination and counter regions, that cell contains $kf(i)$, and
$kf(i)<N_b$. If $s<N_b$ and the state in memory $\mu'$ has every item at its stable
sorting position, every counter at offset $t$ with $s+1<t\le N_b$ equal to the number of
keys below $t$, and every counter at offset $t<s+1$ equal to the number of keys below
$t+1$, then changing the counter at offset $s+1$ to the number of keys below $s+1$ gives
a state in which every item remains at its stable sorting position, counters at offsets
$t>s$ through $N_b$ equal the number of keys below $t$, and counters at offsets $t<s$
equal the number of keys below $t+1$.
\end{theorem}

\begin{theorem}[Final counter shift]
\lean{Light.CsShifted.last}
\label{Light.CsShifted.last}
\leanok
\uses{Light.CountSort.Args.Same, Light.CountSort.Post, Light.CountSort.Pre, Light.CsShifted, Light.CsShifted.placed_update, Light.SameOutside2.update}
\difficulty{4}
Let $\mu$ be an initial memory, $A$ counting-sort arguments, $it,kf$ the item and key
functions, and $\mu\prime$ a current memory. Assume the counting-sort precondition holds,
$\mu$ and $\mu\prime$ agree outside the destination and counter regions, and
$\mu\prime$ satisfies the shifted-state conditions at stage $0$. Then setting the first
counter cell $A.cnt$ to $0$ produces a memory satisfying the counting-sort postcondition:
items are in sorted destination positions, each counter cell contains the number of items
whose key is below its index, and the memory agrees with $\mu$ outside those two regions.
\end{theorem}

\begin{theorem}[Counting-sort shift phase]
\lean{Light.csShift_ends}
\label{Light.csShift_ends}
\leanok
\uses{Light.CountSort.Post, Light.CountSort.Pre, Light.CsPlaced, Light.CsShifted, Light.CsShifted.last, Light.CsShifted.start, Light.CsShifted.succ, Light.CsState, Light.Ends, Light.Ends.next, Light.Ends.storeTo, Light.Ends.whileBlock, Light.SameOutside2.update, Light.csShift, Light.update_frame_setLocal}
\difficulty{5}
Let a counting-sort input have $n$ items with keys in $\{0,\ldots,N-1\}$, and assume the
word and address bounds, key-cell bounds, and required disjointness conditions for the
source, destination, and counter ranges. If the input state has completed placement of
all $n$ items, then the shifting program has an execution of cost at most $18N+9$ whose
final memory has every item at its stable sorted position, has cumulative key counts in
the counter cells, and agrees with the initial memory outside the destination and
counter ranges.
\end{theorem}

\begin{definition}[Counting-sort time bound]
\lean{Light.countSortTime}
\label{Light.countSortTime}
\leanok
For natural numbers $n$ and $	ext{nb}$, the counting-sort time is defined to be
$68n+58\text{nb}+52$.
\end{definition}

\begin{theorem}[Counting sort execution and output]
\lean{Light.countSort_meets}
\label{Light.countSort_meets}
\leanok
\uses{Light.CountSort.Post, Light.CountSort.Pre, Light.CsState, Light.Ends.mono, Light.Ends.next, Light.Meets, Light.Meets.of_body, Light.SameOutside2.refl, Light.countSortBody, Light.countSortTime, Light.csCount_ends, Light.csPlace_ends, Light.csPrefix_ends, Light.csShift_ends, Light.csZero_ends, Light.frame, Light.frame_append_zeros}
\difficulty{4}
Let $P$ be a program, let $q,d$ be natural numbers, and let the machine limits have word
bound $W$, space bound $S$, and depth bound $D$. Let $A$ specify an item count $n$,
source and destination addresses, a key address with stride and offset, a number $b$ of
buckets, and a counter-array address. Suppose the routine at position $q$ in $P$ is the
counting-sort routine; $S\leq W$; the source, destination, and counter intervals of
lengths $n,n,b+1$ lie within the space bound and are pairwise disjoint; and the initial
memory stores item $\mathrm{it}(i)$ at each source address for $i<n$. For each $i<n$,
suppose the key address computed from $\mathrm{it}(i)$ lies within space, stores the
value $kf(i)$, and $kf(i)<b$; moreover, each such key address lies outside both the
destination and counter intervals. Then on the listed arguments and initial memory
$\mu$, the routine has an execution at depth $d$ with cost at most $68n+58b+52$. Its
final memory $\mu'$ places item $\mathrm{it}(j)$ at destination address
$\mathrm{dst}+\mathrm{sortPos}(kf,n,j)$ for every $j<n$, stores at counter address
$\mathrm{cnt}+t$ the number of indices $i<n$ with $kf(i)<t$ for every $t\leq b$, and
agrees with $\mu$ outside the destination and counter intervals.
\end{theorem}
